% GENERATED FILE. Edit canonical lessons, then run npm run generate:books.
% canonical-source-hash: fnv1a64:832b60e9a18cd853
% canonical-lessons: TE-C401-paurudu, TE-C401-prajasvamyam, TE-C401-rajyangam, TE-C401-jatiya, TE-C401-rajakiyalu

\chapter{Citizens and the State}
\label{ch:te-citizens-and-the-state}
\hlchaptermodality{401}

\begin{chapteropening}
\textbf{By the end of this chapter:} \emph{I can talk about citizens, democracy and the constitution.}

Complete the last lesson of chapter 401: I can talk about citizens, democracy and the constitution.
\end{chapteropening}

\section[\texorpdfstring{\te{పౌరుడు} (pauruḍu)}{పౌరుడు (pauruḍu)}]{\te{పౌరుడు} (pauruḍu) — a citizen}

\label{lesson:TE-C401-paurudu}



\begin{quote}
Before the new one: say the Telugu for a campaign, then the Telugu for a majority.
\end{quote}

\subsection*{You'll want to know: \te{పౌరుడు}}
\textbf{\te{పౌరుడు}} — \emph{pauruḍu} — \textquotedblleft{}a citizen\textquotedblright{}.

\begin{quote}
\te{ప్రతి} \te{పౌరుడికి} \te{కొన్ని} \te{హక్కులు} \te{ఉంటాయి}.
\end{quote}

\begin{quote}
\emph{prati pauruḍiki konni hakkulu uṇṭāyi.}
\end{quote}

\begin{quote}
Every citizen has certain rights.
\end{quote}

Plural: \te{పౌరులు}, \emph{paurulu}.

A woman citizen is \te{పౌరురాలు}.

\begin{cousinweb}
\textbf{\te{పౌరుడు}} (\emph{pauruḍu}) is borrowed from Sanskrit \emph{paura}, \textquotedblleft{}a citizen\textquotedblright{} (CDIAL 8419), from \emph{pura}, \textquotedblleft{}a town\textquotedblright{}, with the Telugu masculine ending -\te{డు}.
\end{cousinweb}

\begin{grammarlens}[title={What you've built}]
One more word for giving reasons and opinions.
\end{grammarlens}

\subsection*{Your turn}
\begin{itemize}
\raggedright
  \item \emph{Say it:} \emph{pauruḍu}
  \item \emph{Say it:} \emph{pauruḍu}, once more
  \item \emph{Say it:} \emph{prati pauruḍiki konni hakkulu uṇṭāyi.} — Every citizen has certain rights.
  \item \emph{Recall:} say the Telugu for a campaign, then the Telugu for a majority, then say \emph{pauruḍu} again
\end{itemize}

\begin{tcolorbox}[breakable,colback=teal!4,colframe=teal!35!black,title={Before you move on}]
What does \te{పౌరుడు} mean? (\textquotedblleft{}A citizen\textquotedblright{}.) Say it once more.
\end{tcolorbox}
\section[\texorpdfstring{\te{ప్రజాస్వామ్యం} (prajāsvāmyaṁ)}{ప్రజాస్వామ్యం (prajāsvāmyaṁ)}]{\te{ప్రజాస్వామ్యం} (prajāsvāmyaṁ) — democracy}

\label{lesson:TE-C401-prajasvamyam}



\begin{quote}
Before the new one: say the Telugu for a majority, then the Telugu for a citizen.
\end{quote}

\subsection*{You'll want to know: \te{ప్రజాస్వామ్యం}}
\textbf{\te{ప్రజాస్వామ్యం}} — \emph{prajāsvāmyaṁ} — \textquotedblleft{}democracy\textquotedblright{}.

\begin{quote}
\te{ప్రజాస్వామ్యంలో} \te{జనమే} \te{తమ} \te{నాయకులను} \te{ఎన్నుకుంటారు}.
\end{quote}

\begin{quote}
\emph{prajāsvāmyaṁlō janamē tama nāyakulanu ennukuṇṭāru.}
\end{quote}

\begin{quote}
In a democracy the people themselves choose their leaders.
\end{quote}

\te{ప్రజా}, of the people, + \te{స్వామ్యం}, mastery: rule that belongs to the people.

\begin{cousinweb}
\textbf{\te{ప్రజాస్వామ్యం}} (\emph{prajāsvāmyaṁ}) is built from Sanskrit \emph{prajā}, \textquotedblleft{}offspring, people\textquotedblright{}, and \emph{svāmya}, \textquotedblleft{}mastership, ownership, power\textquotedblright{} (Monier-Williams); the compound is a modern one.
\end{cousinweb}

\begin{grammarlens}[title={What you've built}]
One more word for giving reasons and opinions.
\end{grammarlens}

\subsection*{Your turn}
\begin{itemize}
\raggedright
  \item \emph{Say it:} \emph{prajāsvāmyaṁ}
  \item \emph{Say it:} \emph{prajāsvāmyaṁ}, once more
  \item \emph{Say it:} \emph{prajāsvāmyaṁlō janamē tama nāyakulanu ennukuṇṭāru.} — In a democracy the people themselves choose their leaders.
  \item \emph{Recall:} say the Telugu for a majority, then the Telugu for a citizen, then say \emph{prajāsvāmyaṁ} again
\end{itemize}

\begin{tcolorbox}[breakable,colback=teal!4,colframe=teal!35!black,title={Before you move on}]
What does \te{ప్రజాస్వామ్యం} mean? (\textquotedblleft{}Democracy\textquotedblright{}.) Say it once more.
\end{tcolorbox}
\section[\texorpdfstring{\te{రాజ్యాంగం} (rājyāṅgaṁ)}{రాజ్యాంగం (rājyāṅgaṁ)}]{\te{రాజ్యాంగం} (rājyāṅgaṁ) — a constitution}

\label{lesson:TE-C401-rajyangam}



\begin{quote}
Before the new one: say the Telugu for a citizen, then the Telugu for democracy.
\end{quote}

\subsection*{You'll want to know: \te{రాజ్యాంగం}}
\textbf{\te{రాజ్యాంగం}} — \emph{rājyāṅgaṁ} — \textquotedblleft{}a constitution\textquotedblright{}.

\begin{quote}
\te{మన} \te{రాజ్యాంగం} \te{ప్రతి} \te{పౌరుడికి} \te{స్వేచ్ఛ} \te{ఇస్తుంది}.
\end{quote}

\begin{quote}
\emph{mana rājyāṅgaṁ prati pauruḍiki svēccha istundi.}
\end{quote}

\begin{quote}
Our constitution gives every citizen freedom.
\end{quote}

\te{రాజ్య}, of the state, + \te{అంగం}, a limb, a body.

\begin{cousinweb}
\textbf{\te{రాజ్యాంగం}} (\emph{rājyāṅgaṁ}) is built from Sanskrit \emph{rājya}, \textquotedblleft{}kingship, kingdom\textquotedblright{} (CDIAL 10694), and \emph{aṅga}, \textquotedblleft{}a limb, the body\textquotedblright{} (CDIAL 114): the body of a state.
\end{cousinweb}

\begin{grammarlens}[title={What you've built}]
One more word for giving reasons and opinions.
\end{grammarlens}

\subsection*{Your turn}
\begin{itemize}
\raggedright
  \item \emph{Say it:} \emph{rājyāṅgaṁ}
  \item \emph{Say it:} \emph{rājyāṅgaṁ}, once more
  \item \emph{Say it:} \emph{mana rājyāṅgaṁ prati pauruḍiki svēccha istundi.} — Our constitution gives every citizen freedom.
  \item \emph{Recall:} say the Telugu for a citizen, then the Telugu for democracy, then say \emph{rājyāṅgaṁ} again
\end{itemize}

\begin{tcolorbox}[breakable,colback=teal!4,colframe=teal!35!black,title={Before you move on}]
What does \te{రాజ్యాంగం} mean? (\textquotedblleft{}A constitution\textquotedblright{}.) Say it once more.
\end{tcolorbox}
\section[\texorpdfstring{\te{జాతీయ} (jātīya)}{జాతీయ (jātīya)}]{\te{జాతీయ} (jātīya) — national}

\label{lesson:TE-C401-jatiya}



\begin{quote}
Before the new one: say the Telugu for democracy, then the Telugu for a constitution.
\end{quote}

\subsection*{You'll want to know: \te{జాతీయ}}
\textbf{\te{జాతీయ}} — \emph{jātīya} — \textquotedblleft{}national\textquotedblright{}.

\begin{quote}
\te{జాతీయ} \te{సెలవు} \te{రోజున} \te{బ్యాంకులు} \te{మూసి} \te{ఉంటాయి}.
\end{quote}

\begin{quote}
\emph{jātīya selavu rōjuna byāṅkulu mūsi uṇṭāyi.}
\end{quote}

\begin{quote}
Banks are closed on a national holiday.
\end{quote}

It stands before a noun and does not change: \te{జాతీయ} \te{సెలవు}, a national holiday.

\begin{cousinweb}
\textbf{\te{జాతీయ}} (\emph{jātīya}) is borrowed from Sanskrit \emph{jātīya}, \textquotedblleft{}belonging to a species or kind\textquotedblright{} (Monier-Williams); Telugu uses it for what belongs to the nation.
\end{cousinweb}

\begin{grammarlens}[title={What you've built}]
One more word for giving reasons and opinions.
\end{grammarlens}

\subsection*{Your turn}
\begin{itemize}
\raggedright
  \item \emph{Say it:} \emph{jātīya}
  \item \emph{Say it:} \emph{jātīya}, once more
  \item \emph{Say it:} \emph{jātīya selavu rōjuna byāṅkulu mūsi uṇṭāyi.} — Banks are closed on a national holiday.
  \item \emph{Recall:} say the Telugu for democracy, then the Telugu for a constitution, then say \emph{jātīya} again
\end{itemize}

\begin{tcolorbox}[breakable,colback=teal!4,colframe=teal!35!black,title={Before you move on}]
What does \te{జాతీయ} mean? (\textquotedblleft{}National\textquotedblright{}.) Say it once more.
\end{tcolorbox}
\section[\texorpdfstring{\te{రాజకీయాలు} (rājakīyālu)}{రాజకీయాలు (rājakīyālu)}]{\te{రాజకీయాలు} (rājakīyālu) — politics}

\label{lesson:TE-C401-rajakiyalu}



\begin{quote}
Before the new one: say the Telugu for a constitution, then the Telugu for national.
\end{quote}

\subsection*{You'll want to know: \te{రాజకీయాలు}}
\textbf{\te{రాజకీయాలు}} — \emph{rājakīyālu} — \textquotedblleft{}politics\textquotedblright{}.

\begin{quote}
\te{చాలా} \te{మంది} \te{యువకులకు} \te{రాజకీయాలు} \te{అంటే} \te{ఇష్టం} \te{లేదు}.
\end{quote}

\begin{quote}
\emph{cālā mandi yuvakulaku rājakīyālu aṇṭē iṣṭaṁ lēdu.}
\end{quote}

\begin{quote}
Many young people do not like politics.
\end{quote}

A plural noun. The adjective is \te{రాజకీయ}, political: \te{రాజకీయ} \te{నాయకుడు}, a politician.

\begin{cousinweb}
\textbf{\te{రాజకీయాలు}} (\emph{rājakīyālu}) is built on Sanskrit \emph{rājakīya}, \textquotedblleft{}of or belonging to a king, royal\textquotedblright{} (Monier-Williams); the sense \textquotedblleft{}politics\textquotedblright{} is modern.
\end{cousinweb}

\begin{grammarlens}[title={What you've built}]
One more word for giving reasons and opinions.
\end{grammarlens}

\subsection*{Your turn}
\begin{itemize}
\raggedright
  \item \emph{Say it:} \emph{rājakīyālu}
  \item \emph{Say it:} \emph{rājakīyālu}, once more
  \item \emph{Say it:} \emph{cālā mandi yuvakulaku rājakīyālu aṇṭē iṣṭaṁ lēdu.} — Many young people do not like politics.
  \item \emph{Recall:} say the Telugu for a constitution, then the Telugu for national, then say \emph{rājakīyālu} again
\end{itemize}

\begin{tcolorbox}[breakable,colback=teal!4,colframe=teal!35!black,title={Before you move on}]
What does \te{రాజకీయాలు} mean? (\textquotedblleft{}Politics\textquotedblright{}.) Say it once more.
\end{tcolorbox}
